\section{Checking cocycle connectedness without an orbit search}
\label{sec:cocycle-connectivity}

The native seed stage can prove more from its construction than it could
from arbitrary normal coordinates. Report~31 proves that a minimum-span
surface in a primitive rank-one class is connected. Its spanning-tree
construction also has this property before optimization. We now check the
hypotheses independently and replace the two component searches in this
restricted stage by Euler-characteristic calculations. General normal-surface
queries retain their orbit machinery. The change improves a sufficient
positive test; it does not complete the candidate family.

\subsection{Two certificates excluding zero-class component unions}

Let $M$ be the authenticated compact exterior of the input knot. In
particular it is connected and orientable, with
$H^1(M;\mathbb Z)\cong\mathbb Z$ and $H_2(M;\mathbb Z)=0$.
An integral cocycle $c$ defines the cooriented, edge-coherent local-height
surface $S(c)$ of Section~\ref{sec:cocycle-seeds}. All intersections on
one oriented global edge have the same sign. Deleting components preserves
this property. Signed intersections of any component union form a cocycle,
representing its integral Poincar\'e-dual relative homology class.

First suppose $c$ vanishes on a spanning tree of the one-skeleton. If a
component union $Q$ had zero relative class, its intersection cocycle would
be an integral coboundary $\delta g$. No component meets a tree edge,
since $S(c)$ has zero unsigned intersections there. Thus $\delta g$
vanishes on the tree, forcing $g$ to be constant on all vertices. The
intersection cocycle of $Q$ is zero everywhere. Edge coherence then gives
zero unsigned edge intersections, so $Q$ is empty: every normal piece
meets edges. This argument excludes every nonempty zero-class union
without an optimization witness.

Second suppose a potential $f$ is certified to minimize the total number
of pieces in $S(c+\delta f)$. A nonempty zero-class union would have
cocycle $\delta g$. Deleting it leaves an admissible edge-coherent normal
surface with signed cocycle $c+\delta(f-g)$. Admissible coordinates are
uniquely determined by unsigned edge weights: the kernel of the six
edge-weight equations in a tetrahedron is generated by
$(1,1,1,1,-1,-1,-1)$, and a nonzero displacement would force three
positive quadrilateral types in one of the two vectors. Consequently the
remaining surface is exactly the local-height surface for $f-g$ in normal
coordinates. It has fewer pieces, contradicting optimality. The checked
primal--dual span equality supplies this second exclusion certificate.

In either case each connected component is nonseparating: a separating
cooriented proper surface represents zero relative class. A nonseparating
connected cooriented surface represents a primitive class, because a
closed curve can cross it once and return through its connected complement.
Rank one therefore makes each component a signed generator. Opposite
signs would form a nonempty zero-class union. All signs agree, and a
primitive total class permits exactly one component. Coorientation in
the orientable ambient manifold also proves orientability.

\subsection{An independent integral primitivity check}

The new checker does not invoke cocycle elimination or trust the producer's
reported rank. Rank one follows from the separately authenticated
PD-to-exterior construction. It reconstructs global vertices and signed
global edge equivalences from the finite face records. Every occurrence
of an oriented edge must have the same height difference. Differences
around each local triangle telescope to zero, so these consistency checks
give one integral cocycle; they also guarantee integral constant height
transitions across glued faces.

For a raw witness the graph of zero-valued edges must connect every global
vertex. A breadth-first traversal therefore supplies a zero spanning tree,
independently of the producer's union-find tree. For an optimized witness
the checker instead chooses a tree in the entire edge graph and integrates
the cocycle along its paths. If this integral vertex potential is $p$, the
gauged value of an edge $a\to b$ is
\[
 \widehat c(a,b)=c(a,b)-p(b)+p(a).
\]
The checker requires $\gcd_e\widehat c(e)=1$. Tree-gauged integral
cocycles form the integral kernel of the face equations, a saturated
lattice. Its rank is one, so this gcd proves that the class generates
$H^1(M;\mathbb Z)$. A zero class has gcd zero and is rejected.

Taking the gcd before gauging would be unsound. Doubling a primitive
cocycle and adding a unit vertex coboundary can give raw edge gcd one
while the class remains twice a generator. The regression test constructs
this example on the canonical exterior and rejects both that class and
an exact nonzero coboundary. Local heights may also contain arbitrary
tetrahedron constants. Tests transport raw and optimized proofs through
$20{,}000$-bit offsets and vertex gauges; the checked class and coordinates
remain unchanged, including after hexadecimal JSON serialization.

\subsection{From a connected surface to an essential disc}

The checker validates matching, admissibility and all global edge weights
of the supplied normal vector. For a raw proof, these unsigned weights
must equal $|c(e)|$, which uniquely authenticates its local-height vector.
For an optimized proof, the existing independent span checker reconstructs
the adjusted coordinates, checks its corner matching and potential, and
requires primal--dual equality. Its vector must equal the supplied vector.

Let $F$ be the number of normal pieces, $A$ the sum of their side counts,
and $A_\partial$ the number of normal arcs on unpaired boundary faces.
Then
\[
 \chi=\sum_{e\text{ global}}w_e-\frac{A+A_\partial}{2}+F.
\]
The maintained geometry routine equivalently counts one copy of each
paired face and every boundary face. These are binary sums over incidence
records; no piece expansion or component orbit query is needed.

A connected orientable surface with $\chi=1$ is a disc. Its relative
class is nonzero. Since $H_2(M;\mathbb Z)=0$, the boundary map from
$H_2(M,\partial M;\mathbb Z)$ to $H_1(\partial M;\mathbb Z)$ is
injective, so its boundary is homologically nonzero and hence essential
on the boundary torus. This is a compressing disc of the source knot
exterior and proves that the knot is trivial. The source authentication
is essential to both the rank-one and boundary arguments.

Primitivity and Euler characteristic alone are insufficient. The test
suite starts with the raw cocycle on the canonical exterior of the trefoil
braid $(1,1,1)$ and subtracts two from the potential at global vertex~3.
The resulting coherent surface still represents a primitive class and
has $\chi=1$, but an independent component query finds three components
and no essential disc. Its zero edges do not span the vertices, and it
has no valid optimality witness for those coordinates. Both attempts to
use it as a positive certificate are rejected. This concrete knotted
example guards the connectedness hypothesis itself.

\subsection{Implementation, costs and portfolio behavior}

The new \code{diagram-cocycle-disc-v1} format contains the exact stage-input
PD, canonical triangulation, local heights, normal coordinates and either
a null span witness or the existing primal--dual span certificate. The
independent inspector returns certified connected-surface data for valid
proofs. The public unknot verifier additionally requires $\chi=1$.
It still accepts the previous \code{diagram-normal-disc-v1} format by
replaying its full normal-component proof. New witnesses have intentionally
different hashes; the underlying candidate vectors are preserved.

After geometry authentication, the connectedness check uses $O(T)$
incidence records, breadth-first graph work and gcd operations, plus the
existing span replay's $O(T)$ arithmetic after vertex indexing. For
$B$-bit supplied integers (including heights, potentials and coordinates),
tree potentials and gauged differences have
$O(B+\log(T+1))$ bits. With $\mathsf G(L)$ bounding gcd on $L$-bit
integers, a conservative bound for the added graph, primitivity and Euler
work is
\[
 O\!\left(T\left(B+\log(T+1)+
                   \mathsf G(B+\log(T+1))\right)\right).
\]
This excludes source/manifold authentication and label indexing; it is not
a bound for the whole recognition call. No interval-orbit trace is produced
or replayed in this restricted stage. Cohomology discovery still has a
conservative cubic arithmetic bound, and optimized discovery still uses
the $O(T^2\log(T+1))$ span solver.

The shared work allowance and cancellation callback cover source checks,
both connectedness/Euler replays and optimization. The legacy
\code{max\_cycles} argument is still type-validated; these proofs use
zero interval cycles, so a zero allowance no longer prevents success.
The stage remains disabled by default, tries the raw surface first and
optimizes only on a miss. A non-disc or resource exhaustion yields
\code{INCONCLUSIVE} and returns control to the exact fallback. In
particular the unknot $(1,1,-1)$ still has genus-one raw and optimized
seeds. There is no new general subexponential or quasipolynomial bound.
Section~\ref{sec:cocycle-trees} subsequently adds a bounded family of
alternative tree gauges around this raw-then-span schedule, using the
same independent certificate verifier.

\subsection{Source, adversarial and external controls}

The 22 focused connectivity, seed and span tests pass in 2.165 seconds;
the full suite passes all 1,181 tests in 211.483 seconds. Tests disable
cocycle discovery, flow discovery, exterior construction and both orbit
production and replay while checking serialized positive witnesses. They
also exercise exact schema types, modified height transitions, disconnected
Euler-one surfaces, false optima, divisibility hidden by coboundaries and
cancellation at the first, middle and final guards. Historical component
witnesses, including valid zero-count proofs that must not certify an
unknot, remain covered.

The audit repeats all 84 source queries with the existing two-million-guard
allowance. It obtains 18 positive records, 65 completed restricted misses
and one cap, compared with 17 positives, 64 completed misses and three
caps before this change. The eighteen records contain thirteen distinct
positive witnesses; the named optimized example also occurs as
\code{random-54}, and repeated curls are retained rather than counted
as new constructions. All positives replay independently. The seventeen
previous positives have exactly the same normal vectors, and their
historical proofs are accepted by both old and new verifiers.

The newly successful input is the 32-crossing stabilized circle: its raw
8,385-piece surface now certifies in 53,673 guards. The 8- and 16-crossing
circles use 13,305 and 26,761 guards, compared with 199,317 and 782,613.
The previously capped \code{survivor-00} completes both candidates in
1,295,668 guards, but both remain non-discs. Gordian's raw surface now
finishes with Euler characteristic $-125$; its span optimization still
exhausts the shared allowance. There is no inferred negative verdict and
no newly successful default portfolio case.

Fresh external Regina controls compare both surfaces on every source of at
most twelve crossings: 148 controls from 74 source records. Each agrees
with the new inspector on connectedness, orientability, Euler characteristic
and essential-disc count. These controls supplement, rather than replace,
the independently checked source certificates and the proof of the
connectivity criterion. The final audit takes 36.108 seconds while the
full suite also runs; this elapsed time is not a performance comparison.

On the nineteen ordinary corpus inputs, the previous enabled stage, the
new disabled stage and the new enabled stage give identical expected
verdicts. Every run completes before attempting this stage. The substantial
native-stage gains below therefore do not demonstrate a gain for those
ordinary portfolio calls.

\subsection{Smaller witnesses and complete-call measurements}

The final measurements run serially, without tests, audits, builds or commits
overlapping them. Each case has one warm-up round and five measured rounds.
Each round randomly orders old, old A/A, new and new A/A calls. The tables
give each primary arm's median elapsed time and the median of paired ratios;
the latter need not equal a ratio of the two medians. All 200 measured and
40 warm-up calls complete. Baseline and new certificate hashes are stable
within their own implementations; their coordinate hashes agree across
implementations. Different source-proof schemas intentionally produce
different complete certificate hashes.

The first experiment times source-bound independent verification of a
precomputed positive certificate, together with source parsing, certificate
hashing and serialized-byte counting. Producing the certificates is outside
the timed sample. It includes all geometry and cohomology checks needed by
each respective verifier, with no discovery cache passed into replay.
The mode's total elapsed time, including preparation, is 11.136 seconds.

\begin{center}
\small\begin{tabular}{lrrrrr}
\toprule
Input & Old (ms) & New (ms) & Old/new & Old A/A & New A/A \\
\midrule
circle-8 & 55.735 & 21.468 & 2.591 & 0.993 & 1.006 \\
circle-16 & 121.078 & 43.413 & 2.786 & 0.990 & 0.998 \\
circle-32 & 273.556 & 92.205 & 2.982 & 0.973 & 0.992 \\
Optimized positive & 30.206 & 15.744 & 1.937 & 1.001 & 1.000 \\
\bottomrule
\end{tabular}

\end{center}

Complete replay improves $1.94$--$2.98\times$ on these four cases. The
canonical compact JSON witnesses for the 8-, 16- and 32-crossing circles
shrink from 112,472, 237,105 and 491,293 bytes to 30,077, 60,649 and
121,929 bytes. The optimized positive shrinks from 51,195 to 21,671
bytes, including its height data and dual matching. These sizes include
the common source triangulation and PD; they are not sizes of an isolated
connectivity subcertificate.

The second experiment starts with a source PD and runs the complete
recognizer, including source construction, cohomology discovery, an optional
span search, the stage's internal checks, the exact fallback when required,
and an additional independent replay of any returned native proof. It also
hashes and counts the serialized certificate. Earlier shortcuts are
deliberately disabled so that every case reaches the native stage. Both
arms have no local seed work cap and retain a twelve-second global deadline.
All calls complete before that deadline; the mode takes 53.089 seconds.

\begin{center}
\small\begin{tabular}{lrrrrr}
\toprule
Input & Old (ms) & New (ms) & Old/new & Old A/A & New A/A \\
\midrule
circle-8 & 260.410 & 64.166 & 4.004 & 1.001 & 0.921 \\
circle-16 & 744.828 & 131.303 & 5.756 & 1.002 & 0.998 \\
circle-32 & 2345.215 & 269.486 & 8.702 & 1.005 & 1.002 \\
Optimized positive & 233.967 & 114.397 & 2.050 & 0.999 & 0.999 \\
Genus-one miss & 72.198 & 49.426 & 1.467 & 1.006 & 0.992 \\
Trefoil & 71.814 & 49.088 & 1.459 & 1.001 & 1.008 \\
\bottomrule
\end{tabular}

\end{center}

The four native positives improve $2.05$--$8.70\times$. The 32-crossing
circle's median falls from 2.345 seconds to 0.269 seconds. The genus-one
unknot miss and trefoil still use the exact Khovanov fallback, returning
\code{UNKNOT} and \code{KNOTTED} respectively; their complete calls
improve about $1.46\times$ by making the preceding unsuccessful stage
cheaper. The new-arm A/A ratio on the smallest circle is 0.921, indicating
more variation there than on the other cases; its paired old/new gain is
still about four. These selected-workload gains establish neither ordinary
default-portfolio acceleration nor a general complexity theorem. The
default-budget source audit above separately establishes the newly
completed 32-crossing positive.


Reproduction uses \path{../fast/normal_orbit_research/connectivity.py}, with
\code{audit}, \code{replay} and \code{recognize} modes. The baseline is
commit \code{6212f6f69}; its seed producer and source-proof checker are
loaded from the committed source while all unchanged dependencies remain
fixed. All 478 runtime, driver, fixture and prior-audit pins agree across
the experiments. The incoming archive inventory and the checked
\code{openai/math} head \code{fd4aeeb2e} are unchanged. This integration
uses the existing report~31 theorem and supplies its previously missing
independent source-bound hypothesis checks.

